% ERL Letter: Discussion (implications, limitations including the red team and the H2 pivot,
% generality) and Conclusion are separate files (D-107; revised after the internal review, D-112,
% D-136; near-term framing D-135 addendum 3, D-139 addenda). Every number is a \metric{} macro.

Projections of seabed emergence in a shrinking closed sea have two uncertain inputs, the fall of the water and the shape of the floor it uncovers. The pathway dominates how much seabed emerges by mid-century; the grids decide where. An assessment built on one grid\citep{court2025,duan2025} inherits a location error that a whole-sea area check cannot reveal\citep{bocharov2026}, and a grid used on its own vertical reference can place much of its next metre on seabed that is already dry.

Better estimates need seabed data in the zone that will emerge next, from surveys, charts or lidar bathymetry. The waterline record offers a partial remedy: each metre of fall is surveyed by the next waterlines, as water occurrence and satellite heights have mapped shallow desert lakes\citep{armon2020}, so the audit can be repeated as the sea falls.

The headline rests on the validated coasts because the observations have limits. Errors common to both halves of the observations, from the basin level, persistent wind setup, the threshold or the reed mask, cancel in the floor; a systematic budget for them removes the open shelf from the verdicts but not the west or the interfluve. Surges and setdowns move the North Caspian shoreline by kilometres to tens of kilometres\citep{pavlova2020,yeltay2026}; the registered setup screening could not detect them because its statistic stayed at one grid cell, and our emerged areas assume a static, flat water surface, whereas for infrastructure the surge envelope is the relevant boundary. A split by wind regime, Sentinel-2 in the lowest bins and levelled pre-1992 Landsat scenes would need Earth Engine processing beyond our computing budget and are left to future work. Delta fronts respond to level change by shifting depocentres and prograding\citep{kroonenberg1997,overeem2003}, although the Volga delta plain changed little during the rapid rise that began in 1978\citep{li2004}, and coasts reworked during that rise\citep{kravtsova2012} may differ from the soundings behind the grids; flood-filling a present-day surface cannot predict emergence at the delta fronts.

Our waterline elevation model failed its registered test (H2): against ICESat-2 land heights its NMAD (\metric{h2.atl08.dem.nmad_m}~m) is lower than the grids' (\metric{erl.text.h2.atl08.product_nmad_min_m}--\metric{erl.text.h2.atl08.product_nmad_max_m}~m), but its root-mean-square error is not. H2 is not supported in \metric{rev.g.h2.split.n_not_supported} of \metric{rev.g.h2.split.n_seeds} random reassignments of the test blocks, so failure is the typical outcome, and the NMAD advantage rests on the two northeastern test blocks (Supplementary section~\SIrevision). Emerged seabed may also be lowered after emergence\citep{closson2026}.

Global grids are least constrained where soundings are sparse and vertical references undocumented\citep{weatherall2015,gebco2026}, as in the shallow North Caspian, and a grid that fills recently emerged land from post-emergence elevation data will look skilful above the recent shoreline but not below it. Great Salt Lake shows only that grids without lake-bed data fail where a lidar surface passes, and Lake Urmia, the registered second basin, was not run, so whether the error recurs in other shrinking closed seas is untested.
